%% §9 — re-cut 2026-08-06: the walls, the claim, the program.

\section{Limitations and conclusion}
\label{sec:conclusion}

Three walls stand at every registry size, and the design is drawn up
to them, not past them. \emph{Computability}: a route re-expresses a
question, never changes its degree; unbounded universal claims rest
on certificates, and without one they are bounded claims under
declared bounds, labeled so. \emph{Trust supply}: capability scales
with generated pairs, but the trusted base is the root interpreters,
graded stipulated until corroborated --- and independent lineages,
like external anchors before them, exist in finite supply, so the
grade ladder records where trust saturates rather than pretending it
does not. \emph{Cost}: deciding is exponential in the worst case
everywhere that matters; budgets are permanent provenance, capped is
labeled capped, and a frontier that stops moving under a fixed
registry is a finding --- the measured practice-frontier of that
registry --- not a failure.

Within the walls, the claim is deliberately modest in kind: an
architecture in which one edge kind carries both translation and
solving, one gate admits everything, one currency --- the result ---
holds both progress and failure, and one small proved kernel writes
every entry, so that an autonomous, untrusted LLM can grow the
instrument and play it without any answer resting on its word. The
previous generation of the platform is the first content of the new
one, re-certified through the gate; the standing experiment is
whether the conjecture order --- procedures, then translations, then
languages; semantics first, then syntax --- fed by frontier evidence
alone, reconstructs what was built by hand and then builds what was
not. The loop turns until the player stops cranking; what remains is
the map and the evolved instrument, and the next benchmark starts on
the ground the last one won. The vision is the means. The map is the
end.
